\begingroup
\renewcommand{\arraystretch}{1.45}
\InterpolationArticle{17}
Aus algebraischen Grundsätzen folgt leicht: Bezeichnet man im ersten Fall des vorigen Artikels das Produkt
\[
(\cos t-\cos a)(\cos t-\cos b)(\cos t-\cos c)(\cos t-\cos d)\text{ usw.}
\]
und im zweiten Fall das Produkt
\[
\sin t(\cos t-\cos a)(\cos t-\cos b)(\cos t-\cos c)(\cos t-\cos d)\text{ usw.}
\]
mit \(P\), so muss jede zu \(X\) ähnliche Funktion, die denselben Werten
\(A,B,C,D\) usw. für \(x=a,b,c,d\) usw. genügt, notwendig in der Form \(X+PY\)
enthalten sein, wobei \(Y\) eine unbestimmte Funktion des Bogens \(x\) von der Form
\[
\gamma+\gamma'\cos x+\gamma''\cos 2x+\gamma'''\cos 3x+\text{usw.}
\]
Der umgekehrte Satz, dass nämlich jede in der Form \(X+PY\) enthaltene Funktion jenen Werten genügt, ist selbstverständlich. Ferner ist offenbar das Produkt
\(P\) von \(m+1^{\mathrm{ter}}\) Ordnung, also von der nächsthöheren Ordnung als \(X\); daher gibt es keine andere zu
\(X\) ähnliche Funktion, die denselben Werten genügt, außer einer von höherer Ordnung
als \(X\).

Ist umgekehrt von anderswoher eine Funktion \(X'\) bekannt, die zwar zu \(X\) ähnlich, aber von höherer Ordnung ist
und den vorgegebenen Werten \(A,B,C,D\) usw. genügt, so lässt sich die Funktion \(Y\)
von der von uns angegebenen Form so bestimmen, dass \(X'+PY\) auf
\endgroup
\begingroup
\renewcommand{\arraystretch}{1.45}
\SourceLeaf{302}
\noindent
\(m^{\mathrm{te}}\) Ordnung herabgesetzt wird; daraus muss notwendig eben die Funktion \(X\) hervorgehen, die aus den Werten \(A,B,C,D\) usw.
nach der Methode des vorigen Artikels gewonnen wird. Da sich die Methode zur Bestimmung des höchsten Gliedes
der Funktion \(Y\) und danach der Reihe nach der niedrigeren Glieder als Sonderfall
aus der in Artikel 15, II angegebenen Methode ableiten lässt, ist es nicht nötig, hier
länger dabei zu verweilen.

\Needspace{19\baselineskip}
\InterpolationArticle{18}
Alle vorangegangenen Untersuchungen beruhen auf der Bedingung, dass \(X\) eine Reihe ist, die mit dem Kosinus und Sinus des Bogens \(mx\) abbricht. Gilt diese Bedingung
nicht, so reichen so viele Werte von \(X\), wie wir bisher als bekannt vorausgesetzt haben, offenbar nicht
zur vollständigen Bestimmung aus; der in den verschiedenen von uns betrachteten Fällen angegebene Ausdruck für \(T\) bedarf dann einer Korrektur, die von den folgenden Koeffizienten
in \(X\) abhängt und sich durch die Summationen der Artikel 1 und 2 leicht angeben lässt. Konvergiert dagegen
die Reihe \(X\), ob sie nun ins Unendliche fortläuft oder irgendwo abbricht, so stark,
dass die auf \(\cos mx\) und \(\sin mx\) folgenden Teile als verschwindend angesehen werden können,
so darf auch jene Korrektur vernachlässigt werden, und wir werden damit einen zumindest näherungsweise
richtigen Wert von \(T\) gewonnen haben. Wie es sich damit auch verhalten mag, stellt \(T\) gewiss die einfachste der wahren Funktion ähnliche Funktion
dar, durch die allen vorgegebenen Werten genügt werden kann.

\endgroup
